\documentclass[10pt,a4paper]{amsart}
\usepackage[margin=19mm]{geometry}
\usepackage{amsmath,amssymb,mathtools}
\usepackage[scr=rsfs,cal=euler]{mathalfa}
\usepackage{xfrac}
\usepackage[unicode,hidelinks,bookmarks=false]{hyperref}
\newcommand{\ie}{i.e.\ }
\newcommand{\cf}{cf.\ }
\newcommand{\resp}{respectively\ }
\newcommand{\Subsection}{\subsection}
\providecommand{\nobreakdash}{\nobreak\hbox{-}}
\setlength{\emergencystretch}{2em}
\multlinegap0pt
\usepackage{siunitx}
\usepackage{siunitx}
\usepackage{siunitx}
\usepackage{siunitx}
\usepackage{amssymb}
\usepackage{mathtools}
\usepackage{xcolor}
\usepackage{siunitx}
\usepackage{algorithm}\usepackage{algpseudocode}
\usepackage{float}
\newtheorem{theorem}{Theorem}
\title{Theory of Minimal Weight Perturbations in Deep Networks and its Applications for Low-Rank Activated Backdoor Attacks}
\author{Bethan Evans, Jared Tanner}
\date{Source: arXiv:2601.16880v2 (CC BY 4.0)}
\begin{document}
\maketitle

\noindent\emph{Source and licence.} Theory of Minimal Weight Perturbations in Deep Networks and its Applications for Low-Rank Activated Backdoor Attacks. Version arXiv:2601.16880v2 (CC BY 4.0), distributed by arXiv under the Creative Commons Attribution 4.0 International licence, http://creativecommons.org/licenses/by/4.0/, as recorded in the arXiv licence metadata for this exact version and shown on the abstract page https://arxiv.org/abs/2601.16880v2. Full licence notice: \texttt{licenses/arxiv-2601.16880-CC-BY-4.0.txt}.\par\smallskip

\noindent\textit{Editorial scope.} Counted unit: the article's theorem multi (source0.tex lines 2026--2028). The article's complete original proof of it is delivered with it (source0.tex lines 2029--2033, one \texttt{proof} environment of 5 lines). No mathematics is written, repaired or re-derived: the statement and the proof are the article's own lines, in the article's own notation, and no retained line is substituted or re-flowed.  No further source-local context is needed: every cross-reference of every retained line resolves inside the two delivered spans.  Nothing was omitted from any retained span, so the package carries no omission record.  No retained line contains a citation, so the wrapper carries no bibliography and no bibliography entry.  No retained line contains a figure environment, an image-inclusion command or drawing code, so no image is delivered and the package declares no asset dependency.  Editorial formatting only, in addition: an AMS \texttt{amsart} preamble that restates the article's own declarations and macros used by the retained lines, namely the environments \texttt{theorem} on separate counters, together with the standard packages those lines need.  The retained environments print in this document as Theorem 1 in order of appearance, whereas the article prints the same items with the numbering of its own full document.  The complete native source of the article as distributed in the arXiv e-print for this exact version is bundled unchanged as \texttt{sources/193274/source0.tex}.
\begin{theorem}\label{multi}
For an $M$-layer model $h$ (with, or without activations), let $m \subseteq \{1, \ldots, M\}$ and $n \subseteq m$. Then the minimum total perturbation required to achieve a fixed change in $h(X; \theta)$ by modifying only the layers indexed by $m$ is less than or equal to the minimum total perturbation required when modifying only the layers indexed by $n$.
\end{theorem}

\begin{proof}
 Any feasible perturbation $\{\Delta W_i\}_{i \in n}$ that achieves the desired change in output is also a feasible perturbation for the case where perturbed layers are indexed by $m$, by setting $\Delta W_j = 0$ for all $j \in m \setminus n$.

Therefore, the feasible set for the optimisation over $m$ contains that for $n$, and the optimal value of the perturbation magnitude (e.g., under Frobenius norm) over $m$ is less than or equal to that over $n$. Hence, allowing perturbations in more layers cannot increase the minimum required total perturbation.
\end{proof}

\end{document}
